\documentclass[11pt,a4paper]{article}
\usepackage[utf8]{inputenc}
\usepackage[T1]{fontenc}
\usepackage[brazil]{babel}
\usepackage{lmodern}
\usepackage{microtype}
\usepackage{geometry}
\usepackage{enumitem}
\usepackage{tabularx}
\usepackage{array}
\usepackage{fancyhdr}
\usepackage{titlesec}
\usepackage{xcolor}
\usepackage{hyperref}
\usepackage{xurl}
\usepackage{listings}
\geometry{left=2.3cm,right=2.3cm,top=2.45cm,bottom=2.25cm}
\definecolor{azul}{HTML}{163A5F}
\definecolor{cinza}{HTML}{566573}
\definecolor{fundo}{HTML}{F2F5F8}
\hypersetup{colorlinks=true,urlcolor=azul,linkcolor=azul,pdftitle={Atividade investigativa - Bloqueios de arquivos no Linux},pdfauthor={Sistemas Operacionais}}
\titleformat{\section}{\large\bfseries\color{azul}}{\thesection.}{.6em}{}
\titlespacing*{\section}{0pt}{12pt}{5pt}
\pagestyle{fancy}
\fancyhf{}
\fancyhead[L]{\footnotesize\color{cinza}Sistemas Operacionais \;|\; Atividade investigativa}
\fancyhead[R]{\footnotesize\color{cinza}Linux e bloqueio de arquivos}
\fancyfoot[C]{\footnotesize\thepage}
\renewcommand{\headrulewidth}{.25pt}
\setlength{\headheight}{14pt}
\setlength{\parindent}{0pt}
\setlength{\parskip}{4pt}
\setlist[enumerate]{leftmargin=*,itemsep=3pt,topsep=3pt}
\lstdefinestyle{terminal}{basicstyle=\ttfamily\small,backgroundcolor=\color{fundo},frame=single,rulecolor=\color{fundo},xleftmargin=5pt,xrightmargin=5pt,aboveskip=7pt,belowskip=7pt,showstringspaces=false,breaklines=true}
\begin{document}
\begin{center}
{\Large\bfseries\color{azul} Atividade investigativa}\\[4pt]
{\large\bfseries Bloqueio de arquivos e deadlocks no Linux}\\[8pt]
{\normalsize Comparação experimental entre \texttt{flock()} e \texttt{fcntl()}}\\[7pt]
{\footnotesize Ciência da Computação \;--\; Sistemas Operacionais}
\end{center}
\vspace{2mm}
\color{azul}\hrule height .8pt\color{black}
\vspace{3mm}
\textbf{Nome(s):} \rule{12cm}{.4pt}\\[5pt]
\textbf{Turma:} \rule{5cm}{.4pt}\hfill\textbf{Data:} \rule{4cm}{.4pt}

\section{Objetivo}
Investigar como o Linux trata solicitações de bloqueio sobre arquivos quando duas unidades de execução acessam os mesmos recursos em ordens opostas. Vocês deverão \textbf{formular hipóteses antes de executar}, observar as saídas, estudar as chamadas utilizadas e construir uma explicação para cada resultado.

\textbf{Importante:} esta atividade não apresenta as respostas nem antecipa os resultados dos experimentos. O comportamento observado será parte da investigação.

\section{Entendendo os mecanismos antes do teste}
A chamada \texttt{open()} com a opção \texttt{O\_RDWR} abre um arquivo para leitura e escrita. \textbf{Abrir não é a mesma coisa que adquirir um bloqueio.} Para solicitar bloqueios, os dois programas utilizam mecanismos distintos:

\smallskip
\renewcommand{\arraystretch}{1.23}
\begin{tabularx}{\textwidth}{>{\bfseries}p{2.25cm}X}
\hline
\texttt{flock()} & Solicita bloqueios associados à \textbf{descrição de arquivo aberto}. Com \texttt{LOCK\_EX}, o pedido é de bloqueio exclusivo sobre o arquivo como um todo. No primeiro programa, cada thread faz suas próprias chamadas a \texttt{open()}.\\
\hline
\texttt{fcntl()} & Permite solicitar bloqueios sobre \textbf{regiões de bytes} de um arquivo. A operação tradicional \texttt{F\_SETLKW} com \texttt{F\_WRLCK} solicita bloqueio de escrita e pode aguardar quando há bloqueio conflitante. Esses bloqueios tradicionais são associados ao \textbf{processo}; por isso o segundo programa é executado em dois processos independentes.\\
\hline
\end{tabularx}
\normalsize

Em ambos os casos, estamos estudando bloqueios \textbf{consultivos (advisory)}: processos que ignoram o protocolo de bloqueio ainda podem tentar acessar os arquivos. Nos dois exemplos, os bloqueios são solicitados explicitamente após a abertura.

\section{Preparação}
No terminal do Linux, entrem no diretório \texttt{SO/04\_deadlocks/Atividade} do repositório. Os arquivos \texttt{A.txt} e \texttt{B.txt} já estão preparados. Não é necessário modificar seu conteúdo.
\begin{lstlisting}[style=terminal]
make
\end{lstlisting}
A compilação deve produzir dois programas: \texttt{01\_flock\_threads} e \texttt{02\_fcntl\_processos}. Identifiquem nos códigos onde são utilizados \texttt{open()}, \texttt{sleep(10)}, a solicitação do primeiro bloqueio e a solicitação do segundo.

\newpage
\section{Experimento 1: duas threads e \texttt{flock()}}
\textbf{Arquivo:} \texttt{01\_flock\_threads.c}

Este programa cria duas threads no mesmo processo. A primeira trabalha na ordem \texttt{A.txt} $\rightarrow$ \texttt{B.txt}; a segunda, na ordem \texttt{B.txt} $\rightarrow$ \texttt{A.txt}. Cada thread faz uma pausa de dez segundos depois de solicitar o primeiro bloqueio.

\textbf{Antes da execução, escrevam suas hipóteses:}
\begin{enumerate}[label=H\arabic*.]
\item O que vocês esperam que ocorra após as duas threads solicitarem o segundo bloqueio? Justifiquem.
\item O que significa, neste programa, uma thread manter um bloqueio enquanto solicita outro?
\end{enumerate}

\textbf{Executem:}
\begin{lstlisting}[style=terminal]
make teste_flock
\end{lstlisting}
O comando \texttt{timeout}, usado pelo Makefile, limita externamente o tempo de observação. \textbf{Seu encerramento não deve ser confundido com uma decisão do kernel sobre o bloqueio.}

\textbf{Registrem:} as mensagens exibidas, a ordem dos acontecimentos, se o programa concluiu por conta própria ou se foi interrompido pelo limite de tempo, e se alguma chamada indicou erro.

\textbf{Análise do código:}
\begin{enumerate}[label=\arabic*.,start=1]
\item Por que as threads utilizam \texttt{O\_RDWR} e também \texttt{flock()}? Qual é a diferença entre essas operações?
\item Localizem \texttt{LOCK\_EX}, \texttt{sleep(10)} e \texttt{pthread\_join()}. Expliquem a finalidade de cada um.
\item Identifiquem quais recursos são mantidos e quais são solicitados por cada thread. Desenhem o grafo de espera correspondente \emph{ao estado observado}, se ele existir.
\end{enumerate}

\section{Experimento 2: dois processos e \texttt{fcntl()}}
\textbf{Arquivo:} \texttt{02\_fcntl\_processos.c}

Aqui o mesmo executável é iniciado \textbf{duas vezes}, criando processos independentes. Cada um recebe os nomes dos arquivos por parâmetros, em ordens opostas. Também há uma pausa de dez segundos.

\textbf{Antes da execução, escrevam suas hipóteses:}
\begin{enumerate}[label=H\arabic*.,start=3]
\item O resultado deste experimento será igual ao anterior? Por quê?
\item Por que este experimento executa dois processos em vez de duas threads do mesmo processo?
\end{enumerate}

\textbf{Executem:}
\begin{lstlisting}[style=terminal]
make teste_fcntl
\end{lstlisting}
\textbf{Registrem:} as mensagens de ambos os processos (observando os \emph{PIDs}), se alguma operação aguardou, se apareceu alguma mensagem de erro e o que aconteceu em seguida. Se houver erro, procurem o significado na documentação; não basta copiar a mensagem.

\textbf{Análise do código:}
\begin{enumerate}[label=\arabic*.,start=4]
\item Qual é a finalidade dos campos \texttt{l\_type}, \texttt{l\_whence}, \texttt{l\_start} e \texttt{l\_len} da estrutura \texttt{struct flock}?
\item Qual é a diferença entre \texttt{F\_WRLCK}, \texttt{F\_SETLKW} e a opção \texttt{O\_RDWR}?
\item Em caso de falha na segunda solicitação, quais comandos do programa são executados depois? Expliquem como isso pode afetar o outro processo.
\end{enumerate}

\newpage
\section{Comparação e conclusão}
Respondam às questões a seguir \textbf{com base no código e na saída realmente observada}.
\begin{enumerate}[label=\arabic*.,start=7]
\item Comparem os resultados dos dois experimentos. Ambos concluíram? Algum aguardou, precisou ser interrompido ou retornou erro? O que as evidências indicam sobre cada mecanismo?
\item Retomem as quatro condições necessárias para deadlock (Coffman): exclusão mútua, posse e espera, ausência de preempção e espera circular. Analisem sua presença no(s) estado(s) observado(s).
\item Expliquem por que não é correto afirmar que \texttt{open(..., O\_RDWR)} sozinho impede outro processo ou thread de abrir o arquivo.
\item Consultem a documentação de \texttt{flock()} e de bloqueios tradicionais com \texttt{fcntl()}. Quais diferenças de escopo e de associação dos bloqueios ajudam a interpretar o comportamento observado?
\item Proponham \textbf{uma mudança simples na ordem das solicitações} para reduzir o risco de espera circular. Expliquem por que a mudança funcionaria e, se possível, testem-na em uma cópia local dos arquivos.
\end{enumerate}

\section{Entrega}
Entreguem \textbf{um relatório em PDF} com identificação dos integrantes, hipóteses H1--H4, comandos utilizados, trechos selecionados das saídas, respostas 1--11 e uma conclusão de até um parágrafo sobre o que os testes ensinaram. Não é necessário anexar os códigos inteiros.

Os critérios de avaliação serão: qualidade das hipóteses, fidelidade dos registros, interpretação das chamadas de sistema, compreensão das condições de deadlock e fundamentação da comparação.

\section{Referências para investigação}
\begin{itemize}[leftmargin=*]
\item \textbf{Código-fonte da atividade:} \url{https://github.com/fredluiscabral/UFCIRJ/tree/main/SO/04_deadlocks/Atividade}
\item Manual Linux: \texttt{man 2 open}, \texttt{man 2 flock}, \texttt{man 2 fcntl} e \texttt{man 2 fcntl\_locking}.
\item Documentação online: \url{https://man7.org/linux/man-pages/man2/flock.2.html} e \url{https://man7.org/linux/man-pages/man2/fcntl_locking.2.html}.
\end{itemize}
\vspace{5mm}
\begin{center}
\color{azul}\rule{.88\linewidth}{.4pt}\\[3pt]
\small\textbf{A sequência sugerida é: prever $\rightarrow$ executar $\rightarrow$ registrar $\rightarrow$ explicar.}
\end{center}
\end{document}
